\chapter{Building and Running}
\label{ch:building}

\standfirst{Ten minutes, one dependency, and no network. If \ct{make} does not
work on your machine the reason is in this chapter.}

\section{What you need}

A C11 compiler and \ct{make}. That is the whole hard requirement. SDL3 is
optional---without it the graphics layer builds as a headless stub and
everything except the window still works, which is exactly what the test
suites use. ODBC and OpenSSL are optional in the same way: without the
first there is no database (Chapter~\ref{ch:database}), without the second
no \ct{https} (Chapter~\ref{ch:llm}), and \ct{make deps} reports on all
three at once.

\begin{sh}
git clone https://github.com/blakemcbride/Smalltalk-2026
cd Smalltalk-2026
make
\end{sh}

The last line of a successful build tells you whether you got a window:

\begin{plain}
SDL3: found -- graphics enabled
\end{plain}

If it says \ct{not found -- headless build} and you wanted a window, install
SDL3's development package and build again. On Fedora that is
\ct{sudo dnf install SDL3-devel}; on Debian and Ubuntu, \ct{libsdl3-dev}.

You do \emph{not} need a font, and you do not need Python. The typeface is
rasterised into a checked-in C file; the generator that produced it is only
needed if you want to change the face.

\section{Checking that it works}

\begin{sh}
make test
\end{sh}

This runs two different things and you should know which is which. The
\emph{unit suites} check the system against itself and against the Xerox
traces. The \emph{profile suites} build a real image for each profile and run
that profile's own SUnit tests inside it, holding the score against a
checked-in file. Both must pass.

A clean run ends with every profile at the score
\ct{tests/profiles.expected} records. Note that an \emph{improvement} fails
too, deliberately: a number that got better without being recorded is an
improvement nothing is protecting.

\section{The two object memories}

There is one interpreter and two object memories, chosen at build time. You
will nearly always want the default.

\begin{center}
\small
\begin{tabular}{@{}lp{0.66\textwidth}@{}}
\toprule
\ctp{OM=mt} & 64-bit object table, tagged small integers, per-worker
  allocation, a marking collector, weak references and ephemerons. This is the
  system, and it is the default.\\[3pt]
\ctp{OM=bb} & 16-bit Blue Book memory: object table, reference counting,
  Chapter 30 of the Blue Book verbatim. It exists so the \emph{same
  interpreter} can load the real 1983 Xerox image and reproduce Xerox's
  traces. It cannot bootstrap an image, and says so if you ask it to.\\
\bottomrule
\end{tabular}
\end{center}

\begin{stnote}
Passing \ct{trace2} under \ct{OM=bb} validates the same interpreter, compiler
and BitBlt that the parallel system runs. That is the central trick of the
project, and it is why \ct{OM=bb} is worth keeping around even though nothing
ships it.
\end{stnote}

Other build variables:

\begin{center}
\small
\begin{tabular}{@{}ll@{}}
\toprule
\ctp{TSAN=1} & ThreadSanitizer build\\
\ctp{ASAN=1} & Address and undefined-behaviour sanitizers\\
\ctp{OPT=-O0} & override the optimisation flags\\
\bottomrule
\end{tabular}
\end{center}

Sanitizer builds go to their own directory, so an instrumented binary can
never be linked against stale uninstrumented objects.

\section{Evaluating one expression}

The fastest way to ask the system a question is to build an image in memory,
evaluate something, and exit:

\begin{sh}
./st80 -bootstrap -profile profiles/st2026.profile -eval '3 + 4'
\end{sh}

\begin{plain}
7
\end{plain}

It prints a page of diagnostics on standard error first---how many classes and
methods were loaded, which globals were left undeclared, how many class
initialisers ran. The answer goes to standard output, so \ct{2>/dev/null}
gives you just the answer.

This takes about nine seconds, nearly all of it building the image. There is
no way to keep one warm between invocations, so when you have several
questions, ask them in one program:

Write the program in a file:

\begin{st}
| out |
out := WriteStream on: (String new: 0).
out nextPutAll: (3 + 4) printString.
out nextPutAll: ' and '.
out nextPutAll: 10 factorial printString.
^out contents
\end{st}

\ldots{} and hand it over whole, which sidesteps the shell's opinions about
single quotes:

\begin{sh}
./st80 -bootstrap -profile profiles/st2026.profile -eval "$(cat program.st)"
\end{sh}

\begin{plain}
7 and 3628800
\end{plain}

\begin{stnote}[Sixty-three literals, and no more]
An \ct{-eval} program is compiled as one method, and a method may reference
at most 63 literals---the header holds the count in six bits. Every distinct
string, symbol, large number and literal array in the method is one of them,
and repeated strings are \emph{not} shared, because two occurrences of
\ct{'x'} are two mutable objects. A program that lists fifty things with a
string beside each will hit it and be told so:

\begin{plain}
st80: more than 63 literals referenced; split this method
\end{plain}

Loop over a collection of cases rather than writing them out.
\end{stnote}

An expression containing a caret is treated as a whole method body, temporary
declarations and all. Anything else is a single expression whose value is
wanted, and the system supplies the \ct|^| for you.

\section{Building an image and running the desktop}

\begin{sh}
./st80 -bootstrap -profile profiles/st2026.profile -o st80.image
./st80 -run st80.image
\end{sh}

The first command writes a bootable image---about 11\,MB---and an empty
\ct{st80.image.changes} beside it. The second opens a window on it.

\begin{stwarn}[Which profile?]
\ct{-manifest sources/MANIFEST} builds the pure 1983 system: no exceptions,
no closures, no \ct{Mutex}, and none of the corrections in \ct{lib/}. It is
the museum piece, and \ct{profiles/bluebook.profile} is the supported way to
ask for it. Everything in Parts II, V and VI of this book needs the
\ct{st2026} profile, so build with a profile and not with the manifest.
\end{stwarn}

The window opens on a grey halftone desktop with nothing on it. That is
correct, and Chapter~\ref{ch:tenminutes} is about what to do next.

\begin{sh}
./st80 -run st80.image 2000000
\end{sh}

The optional number is a \emph{bytecode limit}: run this many and stop. It is
for scripted runs and screenshots, not for everyday use.

\begin{stwarn}[The desktop is not on the worker pool yet]
This is the one place where the system's headline claim and what
\ct{-run} does today are not the same thing. Parallel execution is real,
measured and gated---the test suites run at up to 31 native threads under
ThreadSanitizer, and \ct{make bench} sweeps four kernels from one worker to
eight. But the pool is started by those harnesses, not by \ct{-run}: the
windowed desktop runs the interpreter on one thread. Putting the desktop on
the pool is a designed and not-yet-built phase of the project, and
Chapter~\ref{ch:contract} flags what still applies to you today and what is
about the design.
\end{stwarn}

\section{Environment variables for the window}

\begin{center}
\small
\begin{tabular}{@{}lp{0.62\textwidth}@{}}
\toprule
\ctp{ST\_DISPLAY\_THEME} & \ctp{paper} (default), \ctp{classic} (pure black on
  pure white), \ctp{dark}\\
\ctp{ST\_DISPLAY\_SCALE} & screen pixels per display pixel\\
\ctp{ST\_DISPLAY\_WINDOW} & \ctp{WxH}: open the window at exactly this size\\
\ctp{ST\_DISPLAY\_FIT=off} & keep the image's own screen size and letterbox
  it, instead of growing the display to fill the window\\
\ctp{ST\_DISPLAY\_PRESENTATION} & \ctp{integer}, \ctp{letterbox},
  \ctp{stretch}\\
\ctp{ST\_DISPLAY\_TRACE=1} & report why a resize was refused, and print draw
  and present counts at exit\\
\bottomrule
\end{tabular}
\end{center}

\section{Everything \texttt{st80} can do}

\ct{./st80 -help} prints the full list; Appendix~\ref{app:cli} explains each
one. The ones worth knowing on day one:

\begin{center}
\small
\begin{tabular}{@{}ll@{}}
\toprule
\ctp{-bootstrap ... -eval expr} & build an image and evaluate\\
\ctp{-bootstrap ... -o image} & build an image and write it\\
\ctp{-run image [n]} & run an image with \ctp{n} workers\\
\ctp{-syntax f.st...} & compile every method in a file and report failures\\
\ctp{-census image} & summarise an image's object memory\\
\ctp{-classes}, \ctp{-methods} & list what is in an image\\
\ctp{-disasm image Class selector} & dump a method's bytecodes\\
\ctp{-tests} & run every SUnit test and exit non-zero if any failed\\
\ctp{-doctests f.st} & run the \ctp{expr >>> value} examples in a file\\
\bottomrule
\end{tabular}
\end{center}

\section{If the build fails}

\begin{description}
\item[\ctp{fixed object 0 landed on pointer 0}] You asked \ct{OM=bb} to
  bootstrap an image. It cannot. Build with the default \ct{OM=mt}.
\item[Garbled glyphs on screen] A stale object file in the archive. \ct{make clean && make}.
\item[The font and the image disagree] You regenerated the face with
  \ct{make font} and did not rebuild the image. Some classes fix their line
  grid at class-initialisation time, so an image built against one face cannot
  be shown another. Rebuild the image.
\end{description}

Appendix~\ref{app:trouble} has the longer list.
